\subsection{李代数为紧的李群}

定理 1.（H. Weyl）令 G 为李代数是紧半单李代数的连通李群。则 G 紧且其中心有限。

由于 G 半单，其中心 D 离散。此外，商群 G/D 同构于 \(\operatorname{Ad}(G)\)（第III章§6第4节推论4），因而是紧的（命题3）。最后，群 G/D 等于它的换位子群（第III章§9第2节命题4的推论）。于是定理由《积分学》第VII章§3第2节命题5 得出。

命题 4. 令 G 为李代数紧的连通李群。则存在一个环面 T、一个单连通紧半单李群 S、一个向量群 V，以及一个具有有限核的满射态射 \(f: V \times T \times S \to G\)。若 G 紧，则群 V 化为单位元。

令 C（分别地，S）为李代数同构于 L(G) 的中心（分别地，导代数）的单连通李群。则 C 是向量群，S 是中心有限的紧群（定理1），且 G 可与 C × S 关于一个离散子群 D 的商等同起来，D 是中心的（《积分学》第VII章§3第2节引理4）。由于 D 在 S 上投影的像中心因而有限，D∩C 在 D 中的指标有限。令 C' 为 C 中由 D ∩ C 生成的向量子空间，V 为其补子空间。则群 T = C'/(D ∩ C) 是环面，且 G 同构于积群 V × T × S 关于一个有限群的商。

若 G 紧，则 \(V \times T \times S\) 紧（《一般拓扑学》第III章§4第1节命题2的推论2），因而 V 紧，这就推出 \(V = \{e\}\)。

推论 1. 令 G 为连通紧李群。则 \(\mathrm{C}(\mathrm{G})_{0}\) 是环面，\(\mathrm{D}(\mathrm{G})\) 是连通紧半单李群，且态射 \((x,y)\mapsto xy\)（从 \(\mathrm{C}(\mathrm{G})_{0}\times\mathrm{D}(\mathrm{G})\) 到 G 的）是有限覆盖。

沿用命题4 的记号，我们有 \(\mathrm{V} = \{e\}\)，且子群 \(f(\mathrm{T})\) 与 \(f(\mathrm{S})\)（\(\mathrm{G}\) 的）紧，因而闭。于是只需证明 \(f(\mathrm{T}) = \mathrm{C}(\mathrm{G})_0\)、\(f(\mathrm{S}) = \mathrm{D}(\mathrm{G})\)。现在，\(\mathrm{L}(\mathrm{G}) = \mathrm{L}(f(\mathrm{T})) \times \mathrm{L}(f(\mathrm{S}))\)；由于 \(\mathrm{S}\) 半单而 \(\mathrm{T}\) 交换，这就推出 \(\mathrm{L}(f(\mathrm{T})) = \mathcal{C}(\mathrm{L}(\mathrm{G})) = \mathrm{L}(\mathrm{C}(\mathrm{G})_0)\)（第III章§9第3节命题8）且 \(\mathrm{L}(f(\mathrm{S})) = \mathcal{DL}(\mathrm{G}) = \mathrm{L}(\mathrm{D}(\mathrm{G}))\)（第III章§9第2节命题4的推论），由此即得所述断言。

推论 2. 连通紧半单李群的中心与基本群都是有限的。它的泛覆盖是紧的。

沿用命题4 的记号，群 V 与 T 都化为单位元；于是 S 是 G 的泛覆盖，且 G 的基本群同构于 \(\operatorname{Ker} f\)，因而有限。G 的中心 D 由于 G 半单而是离散的，故 D 有限。

推论 3. 连通紧李群 G 的基本群是有限型的 Z-模，其秩等于 C(G) 的维数。

事实上，沿用推论1 的记号，\(\mathrm{C}(\mathrm{G})_{0}\) 的基本群同构于 \(Z^{n}\)，其中 \(n = \dim \mathrm{C}(\mathrm{G})_{0}\)，而 \(\mathrm{D}(\mathrm{G})\) 的基本群有限（推论2）。

推论 4. 令 G 为连通紧李群。下列条件等价：

\begin{enumerate}
\def\labelenumi{(\roman{enumi})}
\item
  G 半单；
\item
  C(G) 有限；
\item
  \(\pi_{1}(G)\) 有限。
\end{enumerate}

若 \(\mathrm{G}\) 单连通，则它半单。

这由推论1 至 3 得出。

推论 5. 令 G 为连通紧李群。则 Int(G) 是李群 Aut(G) 的单位元连通分支（第III章§10第2节）。

令 \(f \in \mathrm{Aut}(\mathrm{G})_0\)。则 \(f\) 诱导出一个自同构 \(f_1\)（\(\mathrm{C}(\mathrm{G})_0\) 的）与一个自同构 \(f_2\)（\(\mathrm{D}(\mathrm{G})\) 的），且我们有 \(f_1 \in \mathrm{Aut}(\mathrm{C}(\mathrm{G})_0)_0\)，\(f_2 \in \mathrm{Aut}(\mathrm{D}(\mathrm{G}))_0\)。由于 \(\mathrm{Aut}(\mathrm{C}(\mathrm{G})_0)\) 离散（《一般拓扑学》第VII章§2第4节命题5），我们有 \(f_1 = \mathrm{Id}\)；由于 \(\mathrm{D}(\mathrm{G})\) 半单，由第III章§10第2节定理1的推论2，存在元素 \(g\)（\(\mathrm{D}(\mathrm{G})\) 的）使 \(f_2(x) = g x g^{-1}\) 对所有 \(x \in \mathrm{D}(\mathrm{G})\) 成立。对所有 \(x \in \mathrm{C}(\mathrm{G})_0\)，我们有 \(g x g^{-1} = x = f_1(x)\)；由于 \(\mathrm{G} = \mathrm{C}(\mathrm{G})_0.\mathrm{D}(\mathrm{G})\)，由此推出 \(g x g^{-1} = f(x)\) 对所有 \(x \in \mathrm{G}\) 成立，即 \(f = \operatorname{Int} g\)。

命题 5. 令 \(\mathbf{G}\) 为李代数紧的李群。

\begin{enumerate}
\def\labelenumi{\alph{enumi})}
\item
  假设 G 连通。则 G 有一个最大紧子群 K；它是连通的。存在 G 的一个闭的中心向量子群（第2节）N，使 G 是直积 \(N \times K\)。
\item
  假设 G 的连通分支群有限。则：
\end{enumerate}

\begin{enumerate}
\def\labelenumi{(\roman{enumi})}
\item
  \(\mathbf{G}\) 的每个紧子群都含于某个极大紧子群。
\item
  若 \(\mathrm{K}_1\) 与 \(\mathrm{K}_2\) 是 \(\mathbf{G}\) 的两个极大紧子群，则存在 \(g \in \mathbf{G}\) 使 \(\mathrm{K}_2 = g\mathrm{K}_1g^{-1}\)。
\item
  令 \(\mathrm{K}\) 为 \(\mathrm{G}\) 的极大紧子群。则 \(\mathrm{K} \cap \mathrm{G}_0\) 等于 \(\mathrm{K}_0\)；它是 \(\mathrm{G}_0\) 的最大紧子群。
\item
  存在 \(G_{0}\) 的一个闭的中心向量子群 N，它在 G 中正规，使得对 G 的任何极大紧子群 K，\(G_{0}\) 是 \(K_{0}\) 与 N 的直积，而 G 是 K 与 N 的半直积。
\end{enumerate}

\begin{enumerate}
\def\labelenumi{\alph{enumi})}
\item
  保留命题4 的记号。\(\mathrm{Ker}f\) 在 V 上的投影是向量群 V 的有限子群，故化为单位元。由此推出 \(\mathrm{Ker}f\) 含于 \(\mathrm{T} \times \mathrm{S}\)，从而 G 是向量群 \(\mathrm{N} = f(\mathrm{V})\) 与紧群 \(\mathrm{K} = f(\mathrm{T} \times \mathrm{S})\) 的直积。G 的每个紧子群在 N 上的投影都化为单位元，故含于 K。这就证明了 \(a\))。
\item
  现在假设 \(G/G_{0}\) 有限。由 a)，\(G_{0}\) 是其最大紧子群 M 与一个向量子群 P 的直积；G 的子群 M 显然正规。令 n 为 \(L(M)\) 在 \(L(G)\) 中的补向量子空间，它在 G 的伴随表示下稳定（第1节及第3节命题3）；它是 \(L(G)\) 的理想，且 \(L(G) = L(M) \times n\)。令 N 为 G 的以 n 为李代数的积分子群；由第III章§6第6节命题14，它在 G 中正规。\(L(G)\) 到 \(L(P)\) 上的以 \(L(M)\) 为核的投影诱导出 n 到 \(L(P)\) 的同构；由此推出 \(G_{0}\) 到 P 上的投影诱导出 N 到 P 的一个 étale 态射；由于 P 单连通，这是一个同构，故 N 是向量群。态射 \((x, y) \mapsto xy\)（从 \(M \times N\) 到 \(G_{0}\) 的）是单的 étale 态射（因为 \(M \cap N\) 化为单位元），因而是同构。由此推出 N 是 G 的闭子群，且商 G/N 紧，因为 \(G_{0}/N\) 紧而 \(G/G_{0}\) 有限（《一般拓扑学》第III章§4第1节命题2的推论2）。
\end{enumerate}

由《积分学》第VII章§3第2节命题3，G 的每个紧子群都含于某个极大紧子群，这些极大紧子群相互共轭，且对 G 的任何极大紧子群 K，G 是 K 与 N 的半直积。由于 \(G_{0}\) 含有 N，它是 N 与 \(G_{0} \cap K\) 的半直积；由此推出 \(G_{0} \cap K\) 连通，从而等于 \(K_{0}\)，因为 \(\mathrm{K}/(\mathrm{G}_{0} \cap \mathrm{K})\) 同构于 \(G/G_{0}\)，因而有限；最后，由 a)，\(K_{0}\) 显然是 \(G_{0}\) 的最大紧子群。

推论. 若 N 满足 b) (iv) 的条件，且 \(K_{1}\) 与 \(K_{2}\) 是 G 的两个极大紧子群，则存在 \(n \in N\) 使 \(nK_{1}n^{-1} = K_{2}\)。

事实上，由 (ii) 存在元素 \(g \in G\) 使 \(gK_{1}g^{-1} = K_{2}\)；由 (iv) 存在 \(n \in N\) 与 \(k \in K_{1}\) 使 g = nk。于是元素 n 具有所要的性质。

